\documentclass[12pt]{article}
\usepackage[T1]{fontenc}
\usepackage[utf8]{inputenc}
\usepackage{lmodern}
\usepackage{textcomp}
\usepackage{enumitem}
\usepackage{geometry}
\geometry{margin=1in}
\begin{document}
\begin{center}
Answer Key - Mathematics - K-12 Level Assessment 23 \
\small (Answers are ordered exactly as in the question paper.)
\end{center}

\section*{Multiple Choice}
\begin{enumerate}[label=\arabic*.]
\item \textbf{Answer:} D
\\ \textit{Options:} A) 15, 17; B) 15, 19; C) 17, 19; D) 17, 21
\\ \textit{Note:} The correct answer is 17, 21 because the pattern increases by 4 with each successive term, and continuing this pattern from 13 results in 17 and then 21.
\item \textbf{Answer:} B
\\ \textit{Options:} A) 1850 x; B) 2000 -150 x; C) 150 x; D) 2000 +150 x
\\ \textit{Note:} The expression 2000 - 150 x correctly represents the total amount of money the Sojourn family has left after spending $150 each day for x days, starting with an initial amount of $2000.
\item \textbf{Answer:} B
\\ \textit{Options:} A) 3 k+1; B) 5 k; C) 3 k; D) 1
\\ \textit{Note:} The correct answer is 5 k because simplifying the expression involves combining like terms, where the fractions are converted to a common denominator and then combined, leading to a final result of 5 k.
\item \textbf{Answer:} B
\\ \textit{Options:} A) add 6 to both sides; B) divide both sides by 6; C) multiply both sides by 6; D) subtract 6 from both sides
\\ \textit{Note:} To isolate the variable a in the equation 6 a = 72, you must perform the inverse operation of multiplication by dividing both sides of the equation by 6.
\item \textbf{Answer:} D
\\ \textit{Options:} A) 11.9 in.; B) 11 in.; C) 45.0 in.; D) 12 in.
\\ \textit{Note:} By rounding each rainfall measurement to the nearest whole number (8 inches, 4 inches, and 0 inches), the estimated total rainfall becomes 12 inches.
\end{enumerate}

\section*{True / False}
\begin{enumerate}[label=\arabic*.]
\setcounter{enumi}{5}
\item \textbf{Answer:} False
\\ \textit{Options:} A) True; B) False
\\ \textit{Note:} The expression 9(9 m + 3 t) simplifies to 81 m + 27 t, not 18 m + 12 t, because when distributing 9, each term inside the parentheses must be multiplied by 9.
\item \textbf{Answer:} True
\\ \textit{Options:} A) True; B) False
\\ \textit{Note:} The statement is true because the cost per stamp is $8.28 ÷ 18 = $0.46, and therefore, for 12 stamps, the total cost is 12 × $0.46 = $5.52.
\item \textbf{Answer:} False
\\ \textit{Options:} A) True; B) False
\\ \textit{Note:} The ratio 12 over 27 simplifies to 4 over 9, while 8 over 18 simplifies to 4 over 9 as well, so they are equivalent ratios, making the statement true, not false.
\item \textbf{Answer:} True
\\ \textit{Options:} A) True; B) False
\\ \textit{Note:} The value of 'a' is -4 because for two lines to be perpendicular, the product of their slopes must equal -1, and the slope of one line is the negative reciprocal of the other, which is satisfied by this value.
\item \textbf{Answer:} False
\\ \textit{Options:} A) True; B) False
\\ \textit{Note:} The total cost of 934 puzzles at $6 each is actually $5,604, calculated by multiplying 934 by 6, which makes the statement false.
\end{enumerate}

\section*{Long Form Response}
\begin{enumerate}[label=\arabic*.]
\setcounter{enumi}{10}
\item \textbf{Answer:} Let $s$ be the side length of the square. Then the dimensions of each rectangle are $s\times\frac{s}{5}$. The perimeter of one of the rectangles is $s+\frac{s}{5}+s+\frac{s}{5}=\frac{12}{5}s$. Setting $\frac{12}{5}s=36$ inches, we find $s=15$ inches. The perimeter of the square is $4 s=4(15\text{ in.})=\boxed{60}$ inches.
\\ \textit{Note:} correct because it accurately derives the side length of the square from the given perimeter of the rectangles, leading to the correct calculation of the square's perimeter.
\item \textbf{Answer:} To find the 8 th number, we start by calculating the total of the first 7 numbers. Since their mean is 15, the total is 7 times 15, which equals 105. When we add the 8 th number, the new mean becomes 12 for all 8 numbers. Therefore, the total for these 8 numbers is 8 times 12, which equals 96. To find the 8 th number, we subtract the total of the first 7 numbers from the total of all 8 numbers: 96 - 105 = -9. Thus, the 8 th number is -9.
\\ \textit{Note:} correct because it accurately calculates the total of the first 7 numbers using the mean, determines the total for all 8 numbers based on the new mean, and then finds the 8 th number by subtracting the total of the first 7 from the total of all 8, resulting in -9.
\end{enumerate}

\vfill
\noindent\textit{End of Answer Key}
\end{document}